\documentclass[11pt]{article}
\usepackage{amsmath,amssymb,amsthm}
\usepackage[utf8]{inputenc}
\newcommand{\dd}{\mathrm{d}}
\newcommand{\Order}{\mathcal{O}}

\newtheorem{proposition}{Proposition}
\newtheorem{corollary}{Corollary}

\title{Asymmetry Between Positive and Negative $K_3$}
\author{HaibiMathAgent}
\date{March 2026}

\begin{document}
\maketitle

\section{Physical Motivation}

Positive $K_3$ represents \emph{cooperative} three-body interaction: groups of
three oscillators with similar phases reinforce each other.  Negative $K_3$
represents \emph{competitive/frustrated} three-body interaction: aligned groups
are penalized.

\section{OA Analysis of the Asymmetry}

From $\dot{r} = -\Delta r + \frac{r(1-r^2)}{2}(K_2 + K_3 r^2)$:

\subsection{$K_3 > 0$ regime}
The effective coupling $K_2 + K_3 r^2$ \emph{increases} with $r$.
This creates a positive feedback loop: more synchronization $\to$ stronger
effective coupling $\to$ even more synchronization.  This is the mechanism
behind explosive synchronization (\Cref{sec:explosive} in main paper).

\paragraph{Fixed point structure:}
\begin{itemize}
  \item For $K_3 < K_2$: single stable FP above threshold (supercritical).
  \item For $K_3 > K_2$: bistability possible (subcritical).
  \item Always: $r^*$ increases with $K_3$ (deeper well).
  \item Always: basin boundary rises with $K_3$ (narrower basin).
\end{itemize}

\subsection{$K_3 < 0$ regime}
The effective coupling $K_2 + K_3 r^2$ \emph{decreases} with $r$.
This creates a \emph{self-limiting} feedback: synchronization weakens
its own coupling, leading to bounded synchronization.

\paragraph{Fixed point structure:}
The fixed-point equation $\Delta = \frac{(1-r^2)}{2}(K_2 + K_3 r^2)$
has qualitatively different solutions:

\begin{proposition}[Self-limiting synchronization for $K_3<0$]
  For $K_3 < 0$:
  \begin{enumerate}
    \item The synchronized fixed point $r^*$ \emph{decreases} with $|K_3|$
      (weaker synchronization).
    \item The basin of attraction of the synchronized state \emph{expands}
      with $|K_3|$ (easier to reach).
    \item The bifurcation is always supercritical (never explosive).
    \item There exists a critical $K_3^{\min} = -K_2/(r^{*2})$ below which
      the effective coupling becomes repulsive at $r^*$, and the synchronized
      state is destroyed.
  \end{enumerate}
\end{proposition}

\begin{proof}
  (1) From $\frac{\dd r^*}{\dd K_3} = \frac{r^{*2}(1-r^{*2})}{2D}$ with
  $D>0$ in the supercritical regime, and noting $r^{*2}>0$, $1-r^{*2}>0$:
  $\dd r^*/\dd K_3 > 0$, so $r^*$ decreases as $K_3$ becomes more negative.

  (2) The basin boundary $r_{\mathrm{sep}}$ satisfies the same equation
  but at an unstable FP.  For $K_3<0$, the cubic structure
  $K_3 u^2 - (K_2+K_3)u + (K_2-2\Delta)=0$ with $K_3<0$ has the
  leading coefficient negative, pushing $r_{\mathrm{sep}}$ toward 0
  (wider basin).

  (3) The cubic coefficient in the normal form is $-(K_2-K_3)/2$.
  For $K_3<0$: $K_2-K_3 = K_2+|K_3| > 0$ always, so the bifurcation
  is supercritical.

  (4) At $K_2 + K_3 r^{*2} = 0$, i.e., $K_3 = -K_2/r^{*2}$, the
  effective coupling vanishes. Below this, the synchronized state
  cannot be sustained.
\end{proof}

\section{Physical Interpretation: ``Deeper Narrower'' vs ``Shallower Wider''}

The two regimes represent fundamentally different synchronization strategies:

\begin{center}
\begin{tabular}{lll}
  & $K_3 > 0$ & $K_3 < 0$ \\
\hline
  $r^*$ & Higher (stronger sync) & Lower (weaker sync) \\
  Basin & Narrower (harder to reach) & Wider (easier to reach) \\
  Stability & Deeper well & Shallower well \\
  Transition & Can be explosive & Always continuous \\
  Strategy & ``Elite synchronization'' & ``Democratic synchronization'' \\
\end{tabular}
\end{center}

\paragraph{Implications for neural systems:}
Negative $K_3$ (competitive higher-order) may be preferred in biological
systems that need \emph{robust} synchronization (wide basin, easy to enter)
over \emph{strong} synchronization (high $r^*$).  This could explain why
neural circuits exhibit moderate, fluctuating synchrony rather than full
phase locking.

\section{Cluster States for $K_3 < 0$}

GPU-Claude-Opus (Task M.1) found multi-cluster states (2--3 clusters) only
at moderate $K_2$ + negative $K_3$.  On the OA manifold, cluster states are
not captured (OA assumes identical oscillator populations).  The emergence
of clusters for $K_3<0$ signals departure from the OA manifold, consistent
with the $r_2/r_1^2 > 1$ diagnostic.

A full cluster theory requires extending the OA framework to multiple
populations, which is beyond the scope of this paper but represents an
important future direction.

\end{document}
